% !TeX program = pdflatex
\documentclass[11pt,a4paper]{article}
\usepackage[utf8]{inputenc}
\usepackage[T1]{fontenc}
\usepackage{amsmath,amssymb,amsthm}
\usepackage[margin=2.4cm]{geometry}
\usepackage{booktabs}
\usepackage{array}
\usepackage{enumitem}
\usepackage{xcolor}
\usepackage[hidelinks,breaklinks]{hyperref}
\usepackage[expansion=false]{microtype}

\sloppy
\emergencystretch=2em

\newtheorem{definicao}{Definição}
\newtheorem{teorema}{Teorema}
\newtheorem{proposicao}{Proposição}
\newtheorem{lema}{Lema}
\newtheorem{observacao}{Observação}
\newtheorem{corolario}{Corolário}
\newtheorem{questao}{Questão}

\newcommand{\abs}[1]{\lvert #1\rvert}

\title{Álgebra do Suporte Coordenado Recursivo}
\author{}
\date{}

\begin{document}
\maketitle

\begin{abstract}
\noindent
Formalizamos uma família recursiva de suportes coordenados construída por
aninhamento de pares sobre um alfabeto finito arbitrário.
A construção separa explicitamente alfabeto, codificação, estrutura recursiva,
coordenada e representação.
Sobre esta família definem-se a álgebra de campos (com valores num corpo ou
numa álgebra associativa com unidade) e a álgebra de valores obtida por
extensões quadráticas.
Demonstramos a bijecção da coordenada posicional, o escalonamento sob a
inclusão canónica de peso mínimo, a estrutura de álgebra de campos com base
de idempotentes, a decomposição da álgebra de valores quando o discriminante
é um quadrado, e a existência do colimite de conjuntos do sistema indutivo
das inclusões.
Distinguimos a inclusão de peso mínimo, que multiplica a coordenada pela
cardinalidade do alfabeto, da inclusão alternativa de peso máximo, que
preserva a coordenada.
Questões em aberto incluem a existência de uma coordenada canónica no
colimite, a extensão da categoria restrita dos níveis da família, e a
passagem da soma parcial algébrica a uma exponencial sem topologia adicional.
\end{abstract}

\tableofcontents
\newpage

\section{Introdução}

Um \emph{suporte coordenado} é um conjunto finito munido de uma bijecção para
um intervalo inicial de inteiros.
Essa estrutura mínima permite indexar estados discretos e definir campos
(funções com valores num corpo ou numa álgebra) sem impor, a priori, produto
cartesiano plano, ordem intrínseca ou estrutura linear no conjunto de estados.

Neste artigo construímos uma família recursiva de suportes coordenados por
aninhamento de pares a partir de um alfabeto finito.
A recursão acrescenta um símbolo na posição de peso mínimo, de modo que a
coordenada do nível anterior é multiplicada pela cardinalidade do alfabeto.
Sobre cada nível definem-se o espaço vectorial dos campos escalares e a
álgebra de campos com valores numa álgebra associativa com unidade.
Uma segunda camada, independente do suporte, produz álgebras de valores por
quocientes quadráticos e, sob hipóteses adequadas, uma duplicação do índice
do suporte.

A contribuição principal é a formalização autocontida desta álgebra:
definições, demonstrações e distinção rigorosa entre o que é estrutura do
suporte, o que é coordenada, o que é campo e o que é valor.
Em particular, separamos a inclusão de peso mínimo (escalonamento da
coordenada) da inclusão alternativa de peso máximo (preservação da
coordenada), e registamos que o colimite de conjuntos do sistema indutivo
não herda automaticamente uma coordenada compatível com todas as aplicações
coordenadas dos níveis finitos.

\paragraph{Pré-requisitos.}
Assumimos familiaridade com anéis comutativos com unidade, ideais, anéis
quociente, corpos, e álgebras associativas com unidade sobre um corpo.
O Teorema Chinês dos Restos é enunciado na forma utilizada
(Secção~\ref{sec:valores}).

\paragraph{Organização.}
A Secção~\ref{sec:axiomas} fixa o suporte coordenado e os morfismos.
A Secção~\ref{sec:alfabeto} introduz alfabeto, codificação e par-base.
A Secção~\ref{sec:familia} constrói a família recursiva e a inclusão canónica.
A Secção~\ref{sec:campos} define a álgebra de campos.
A Secção~\ref{sec:valores} trata das álgebras de valores e da decomposição
quadrática.
A Secção~\ref{sec:categoria} discute o sistema indutivo, o colimite e a
categoria restrita dos níveis.
A Secção~\ref{sec:realizacoes} introduz realizações e campos de contagem.
A Secção~\ref{sec:abertas} recolhe questões em aberto.

% ═══════════════════════════════════════════════════════════════════════════
\section{Axiomas mínimos e objecto fundamental}
\label{sec:axiomas}

\begin{definicao}[Suporte coordenado]
\label{def:suporte-coordenado}
Um \emph{suporte coordenado} é um par \((X,C)\) em que \(X\) é um conjunto
finito não-vazio e
\[
C\colon X\overset{\sim}{\longrightarrow}\{0,1,\dots,\lvert X\rvert-1\}
\]
é uma bijecção.
\end{definicao}

\begin{observacao}[Minimalidade]
\label{obs:minimalidade}
A Definição~\ref{def:suporte-coordenado} exige apenas finitude e a existência
de uma bijecção para um intervalo inicial de inteiros.
Nenhuma estrutura de ordem, produto, inclusão ou operação algébrica é
imposta ao conjunto \(X\) neste nível.
\end{observacao}

\begin{definicao}[Morfismos de suportes coordenados]
\label{def:morfismos}
Sejam \((X,C)\) e \((Y,D)\) suportes coordenados e seja \(q\ge 2\) um inteiro
fixado no contexto.
\begin{enumerate}[leftmargin=1.5em]
\item Um \emph{morfismo estrito} é uma função \(f\colon X\to Y\) tal que
\[
D\circ f = C.
\]
Como \(D\) é bijecção, esta igualdade equivale a
\(f = D^{-1}\circ C\), e a condição de existência reduz-se a
\(\operatorname{im} C\subseteq\operatorname{im} D\)
(o que é automático quando os contradomínios coincidem).

\item Um \emph{morfismo de escala \(k\)} (\(k\in\mathbb N\), \(k\ge 1\))
é uma função \(f\colon X\to Y\) para a qual se verificam simultaneamente:
\begin{enumerate}[label=(\alph*)]
\item a condição de imagem
\[
q^k\cdot\operatorname{im} C
\subseteq
\operatorname{im} D;
\]
\item a igualdade
\[
D\circ f = q^k\cdot C.
\]
\end{enumerate}
Sob a condição~(a), a igualdade~(b) determina \(f\) de forma única:
\[
f = D^{-1}\circ(q^k\cdot C).
\]
\end{enumerate}
As duas classes são disjuntas quando \(k\ge 1\) e \(\lvert X\rvert\ge 2\):
nesse caso \(q^k\cdot C\neq C\) (pois \(C\) não é identicamente nula e
\(q^k\neq 1\)), logo nenhuma função pode satisfazer simultaneamente
\(D\circ f=C\) e \(D\circ f=q^k\cdot C\).
Se \(\lvert X\rvert=1\), então \(C\) é a função constante \(0\) e
\(q^k\cdot C=C\), pelo que a distinção entre as duas classes colapsa.
\end{definicao}

\begin{observacao}[Existência e unicidade]
\label{obs:compatibilidade}
A condição de imagem \(q^k\cdot\operatorname{im} C\subseteq\operatorname{im} D\)
é necessária e suficiente para a existência de \(f\).
Quando ela se verifica, \(f\) é única porque \(D\) é injectiva.
A notação \(q^k\cdot\operatorname{im} C\) designa o conjunto
\(\{q^k\cdot c:c\in\operatorname{im} C\}\subseteq\mathbb N\), isto é, a
imagem de \(C\) sob a multiplicação escalar de inteiros; não se trata de
uma operação sobre o suporte \(X\).
\end{observacao}

% ═══════════════════════════════════════════════════════════════════════════
\section{Alfabeto, codificação e par-base}
\label{sec:alfabeto}

\begin{definicao}[Alfabeto e codificação]
\label{def:alfabeto}
Seja \(B\) um conjunto finito com
\[
q := \lvert B\rvert \ge 2.
\]
Fixe uma bijecção (codificação dos símbolos)
\[
d\colon B\overset{\sim}{\longrightarrow}\{0,1,\dots,q-1\}.
\]
\end{definicao}

\begin{observacao}[Independência]
\label{obs:independencia-alfabeto}
O alfabeto \(B\) e a codificação \(d\) são escolhas independentes da
estrutura recursiva que será construída a seguir.
Casos particulares (não privilegiados):
\begin{itemize}[leftmargin=1.5em]
\item binário: \(q=2\), \(B=\{0,1\}\), \(d=\mathrm{id}\);
\item hexadecimal: \(q=16\);
\item byte: \(q=256\), \(B=\{0,1,\dots,255\}\), \(d=\mathrm{id}_B\);
equivalentemente, \(B=\mathbb Z/256\mathbb Z\), com \(d\) dado pelo
representante canónico de cada classe em \(\{0,\dots,255\}\).
\end{itemize}
\end{observacao}

\begin{definicao}[Par-base]
\label{def:par-base}
\[
X^{(2)} := B\times B,
\qquad
C^{(2)}(b_1,b_2) := d(b_1) + q\, d(b_2).
\]
O par \(\bigl(X^{(2)},C^{(2)}\bigr)\) é um suporte coordenado.
\end{definicao}

\begin{proposicao}[Bijecção do par-base]
\label{prop:bijecao-base}
A aplicação \(C^{(2)}\) é bijecção sobre \(\{0,1,\dots,q^2-1\}\).
\end{proposicao}

\begin{proof}
Trata-se da representação posicional clássica em base \(q\) com dois
dígitos.
Para cada \(c\in\{0,\dots,q^2-1\}\) existem únicos \(r_0,r_1\in\{0,\dots,q-1\}\)
com \(c=r_0+q\,r_1\).
Como \(d\) é bijecção, existem únicos \(b_1,b_2\in B\) com
\(d(b_1)=r_0\) e \(d(b_2)=r_1\).
Logo \(C^{(2)}\) é bijectiva.
\end{proof}

% ═══════════════════════════════════════════════════════════════════════════
\section{Família recursiva e inclusão canónica}
\label{sec:familia}

\begin{definicao}[Família recursiva por aninhamento]
\label{def:familia}
Definimos indutivamente \(X^{(n)}\) e \(C^{(n)}\) para \(n\ge 2\).

\begin{enumerate}[leftmargin=1.5em]
\item Base (\(n=2\)): conforme Definição~\ref{def:par-base}.

\item Passo (\(n\ge 3\)):
\[
X^{(n)} := B\times X^{(n-1)}.
\]
Um elemento escreve-se \(x=(b,\xi)\) com \(b\in B\) e \(\xi\in X^{(n-1)}\).
A coordenada é
\[
C^{(n)}(b,\xi) := d(b) + q\, C^{(n-1)}(\xi).
\]
\end{enumerate}
\end{definicao}

\begin{observacao}[Notação de aninhamento]
\label{obs:aninhamento}
Um estado de comprimento \(n\) escreve-se
\[
x = \bigl(b_1,\bigl(b_2,\bigl(b_3,\dots,(b_{n-1},b_n)\dots\bigr)\bigr)\bigr),
\]
com \(b_i\in B\).
Aqui \(b_1\) é o símbolo de peso \(q^0\) (menos significativo) e \(b_n\) o de
peso \(q^{n-1}\) (mais significativo).
A coordenada correspondente é
\[
C^{(n)}(x) = \sum_{i=1}^{n} q^{i-1}\, d(b_i).
\]
A estrutura sintáctica do aninhamento e a expansão posicional estão alinhadas.
\end{observacao}

\begin{proposicao}[Bijecção]
\label{prop:bijecao}
Para todo \(n\ge 2\),
\[
C^{(n)}\colon X^{(n)}\overset{\sim}{\longrightarrow}\{0,1,\dots,q^n-1\}
\]
é bijecção.
Em particular \(\lvert X^{(n)}\rvert = q^n\).
\end{proposicao}

\begin{proof}
Por indução em \(n\).

\begin{itemize}[leftmargin=1.5em]
\item Base (\(n=2\)): Proposição~\ref{prop:bijecao-base}.

\item Passo: assuma o resultado no nível \(n-1\) (\(n\ge 3\)), isto é,
que \(C^{(n-1)}\) é bijecção sobre \(\{0,\dots,q^{n-1}-1\}\).
A aplicação
\[
C^{(n)}\colon(b,\xi)\mapsto d(b)+q\, C^{(n-1)}(\xi)
\]
é a representação posicional clássica em base \(q\): \(d(b)\) ocupa o dígito
de peso \(q^0\) e \(C^{(n-1)}(\xi)\) ocupa os dígitos de pesos
\(q^1,\dots,q^{n-1}\).
Como \(d(b)\) percorre \(\{0,\dots,q-1\}\) e \(C^{(n-1)}(\xi)\) percorre um
bloco completo de comprimento \(q^{n-1}\), a imagem é
\(\{0,\dots,q^{n}-1\}\) e a aplicação é injectiva, logo bijectiva.
\end{itemize}
\end{proof}

\begin{definicao}[Inclusão canónica de peso mínimo]
\label{def:iota}
Para cada \(n\ge 2\), seja \(b_*=d^{-1}(0)\) o (único) símbolo de código zero.
A inclusão canónica é
\[
\iota_n\colon X^{(n)}\hookrightarrow X^{(n+1)},
\qquad
\iota_n(\xi) := (b_*,\xi).
\]
A inclusão é determinada unicamente pela codificação \(d\).
\end{definicao}

\begin{proposicao}[Escalonamento da coordenada]
\label{prop:escalonamento}
Para todo \(n\ge 2\),
\[
C^{(n+1)}\circ\iota_n = q\cdot C^{(n)}.
\]
\end{proposicao}

\begin{proof}
Por definição,
\begin{align*}
C^{(n+1)}\bigl(\iota_n(\xi)\bigr)
&= C^{(n+1)}(b_*,\xi)
= d(b_*) + q\, C^{(n)}(\xi)
= 0 + q\, C^{(n)}(\xi)
= q\, C^{(n)}(\xi).
\end{align*}
\end{proof}

\begin{corolario}[Coerência da família]
\label{cor:coerencia}
Para todo \(k\ge 1\),
\[
C^{(n+k)}\circ\iota_{n+k-1}\circ\cdots\circ\iota_n = q^k\cdot C^{(n)}.
\]
\end{corolario}

\begin{proof}
Indução em \(k\). O caso \(k=1\) é a Proposição~\ref{prop:escalonamento}.
O passo indutivo segue por composição.
\end{proof}

\begin{proposicao}[Ordem induzida]
\label{prop:ordem}
A relação
\[
\xi\le_{C^{(n)}}\eta
\quad\Longleftrightarrow\quad
C^{(n)}(\xi)\le C^{(n)}(\eta)
\]
é uma ordem total em \(X^{(n)}\).
Sob a inclusão,
\[
\xi\le_{C^{(n)}}\eta
\quad\Longrightarrow\quad
\iota_n(\xi)\le_{C^{(n+1)}}\iota_n(\eta),
\]
porque \(C^{(n+1)}\circ\iota_n=q\cdot C^{(n)}\) e a multiplicação por \(q>0\)
preserva a ordem.
\end{proposicao}

\begin{proof}
A totalidade segue de \(C^{(n)}\) ser bijecção sobre um intervalo de inteiros.
A preservação da ordem sob \(\iota_n\) é consequência directa da
Proposição~\ref{prop:escalonamento}.
\end{proof}

\begin{observacao}[Inclusões como morfismos de escala]
\label{obs:iota-escala}
Pela Proposição~\ref{prop:escalonamento} e pela condição de imagem
\(q\cdot\operatorname{im} C^{(n)}\subseteq\operatorname{im} C^{(n+1)}\)
(satisfeita porque
\(\operatorname{im} C^{(n)}=\{0,\dots,q^n-1\}\) e
\(\operatorname{im} C^{(n+1)}=\{0,\dots,q^{n+1}-1\}\)),
cada \(\iota_n\) é um morfismo de escala \(1\) no sentido da
Definição~\ref{def:morfismos}.
\end{observacao}

\begin{definicao}[Inclusão alternativa de peso máximo]
\label{def:iota-peso-max}
Fixe o mesmo alfabeto \(B\) e a mesma codificação \(d\).
Para cada \(n\ge 2\), seja \(Y^{(n)}\) o conjunto das \(n\)-tuplas
\((b_0,\dots,b_{n-1})\) com \(b_i\in B\), equipado com a coordenada
\[
C'_n(b_0,\dots,b_{n-1})=\sum_{k=0}^{n-1}d(b_k)\,q^k.
\]
A \emph{inclusão de peso máximo} é
\[
\iota'_n\colon Y^{(n)}\hookrightarrow Y^{(n+1)},
\qquad
\iota'_n(b_0,\dots,b_{n-1})=(b_0,\dots,b_{n-1},b_*),
\]
onde \(b_*=d^{-1}(0)\).
\end{definicao}

\begin{proposicao}[Preservação da coordenada sob inclusão de peso máximo]
\label{prop:peso-max}
Com as notações da Definição~\ref{def:iota-peso-max},
\[
C'_{n+1}\circ\iota'_n = C'_n.
\]
\end{proposicao}

\begin{proof}
O símbolo acrescentado ocupa a posição de peso \(q^n\) e tem código zero, logo
não altera a soma posicional dos \(n\) símbolos anteriores.
\end{proof}

\begin{observacao}[Duas famílias distintas]
\label{obs:comparacao-familias}
A inclusão de peso mínimo (Definição~\ref{def:iota}) satisfaz
\(C^{(n+1)}\circ\iota_n=q\cdot C^{(n)}\).
A inclusão de peso máximo (Definição~\ref{def:iota-peso-max}) satisfaz
\(C'_{n+1}\circ\iota'_n=C'_n\).
Os conjuntos \(X^{(n)}\) (aninhados) e \(Y^{(n)}\) (tuplas planas) não são
literalmente iguais; são canonicamente identificáveis pela aplicação de
achatamento
\[
(b_1,(b_2,\dots,(b_{n-1},b_n)\dots))\longmapsto(b_1,\dots,b_n),
\]
que preserva a coordenada posicional.
Sob essa identificação, na instanciação \(q=256\), \(d\) como abaixo, as leis
posicionais coincidem nível a nível (para \(n\ge 2\)), mas os sistemas de
transição permanecem distintos.
Ambas as construções produzem famílias de suportes coordenados no sentido da
Definição~\ref{def:suporte-coordenado}; a escolha da inclusão é estrutura
adicional.
\end{observacao}

\begin{proposicao}[Instanciação particular]
\label{prop:instanciacao-bytes}
Seja \(q=256\) e seja \(B=\{0,1,\dots,255\}\) (o conjunto de representantes
canónicos módulo \(256\)), com \(d=\mathrm{id}_B\).
Equivalentemente, pode-se tomar \(B=\mathbb Z/256\mathbb Z\) e definir
\(d([r])\) como o único representante de \([r]\) em \(\{0,\dots,255\}\).
Em ambos os casos obtém-se uma família de suportes coordenados: para cada
\(n\ge 2\), o par \(\bigl(X^{(n)},C^{(n)}\bigr)\) satisfaz a
Definição~\ref{def:suporte-coordenado}.
\end{proposicao}

\begin{proof}
Consequência directa da Proposição~\ref{prop:bijecao}, uma vez que \(d\) é
bijecção sobre \(\{0,\dots,q-1\}\).
\end{proof}

\begin{observacao}[Camadas distintas]
\label{obs:camadas}
A construção distingue explicitamente:
\begin{enumerate}[leftmargin=1.5em]
\item alfabeto \(B\);
\item codificação \(d\);
\item estrutura recursiva \(X^{(n)}\);
\item coordenada \(C^{(n)}\);
\item representação \((b_1,\dots,b_n)\).
\end{enumerate}
O suporte não é identificado com a representação numérica.
A base estrutural é o par \(X^{(2)}=B\times B\); não existe nível de
comprimento \(1\) na família.
\end{observacao}

% ═══════════════════════════════════════════════════════════════════════════
\section{Álgebra de campos}
\label{sec:campos}

\begin{definicao}[Espaço de campos escalares]
\label{def:espaco-campos}
Seja \(\mathbb K\) um corpo e seja \(X^{(n)}\) um suporte da família.
O espaço de campos escalares sobre \(X^{(n)}\) é
\[
\mathcal V_n := \mathbb K^{X^{(n)}}
= \{V\mid V\colon X^{(n)}\to\mathbb K\}.
\]
As operações são ponto a ponto:
\[
(V+W)(x) := V(x)+W(x),
\qquad
(\lambda V)(x) := \lambda\, V(x).
\]
Assim \(\mathcal V_n\) é um espaço vectorial sobre \(\mathbb K\).
\end{definicao}

\begin{observacao}[Separação suporte/campo]
\label{obs:suporte-campo}
Tem-se \(X^{(n)}\neq\mathcal V_n\).
Nenhuma estrutura linear é atribuída ao suporte apenas pela definição de
\(\mathcal V_n\).
\end{observacao}

\begin{proposicao}[Dimensão e base canónica]
\label{prop:dimensao}
Para todo \(n\ge 2\),
\[
\dim_{\mathbb K}\mathcal V_n = \lvert X^{(n)}\rvert = q^n.
\]
A família \(\{e_x\mid x\in X^{(n)}\}\) definida por
\[
e_x(y) :=
\begin{cases}
1 & \text{se }y=x,\\
0 & \text{se }y\neq x
\end{cases}
\]
é uma base de \(\mathcal V_n\).
\end{proposicao}

\begin{proof}
Como \(X^{(n)}\) é finito, o espaço de funções \(\mathbb K^{X^{(n)}}\)
possui a base canónica indexada pelos elementos de \(X^{(n)}\).
A verificação de que \(\{e_x\}\) é linearmente independente e gera é
imediata.
\end{proof}

\begin{definicao}[Álgebra de campos]
\label{def:algebra-campos}
Seja \(Y\) uma \(\mathbb K\)-álgebra associativa com unidade.
A álgebra de campos com valores em \(Y\) é
\[
\mathcal A_n(Y) := Y^{X^{(n)}}
= \{F\mid F\colon X^{(n)}\to Y\},
\]
com adição e multiplicação ponto a ponto e unidade o campo constante
\(\mathbf 1(x)=1_Y\).
\end{definicao}

\begin{proposicao}[Idempotentes canónicos]
\label{prop:idempotentes}
Os elementos \(e_x\) da base de \(\mathcal V_n\) (Proposição~\ref{prop:dimensao}),
quando vistos como elementos da álgebra \(\mathcal A_n(\mathbb K)\),
satisfazem
\[
e_x^2 = e_x,
\qquad
e_x\, e_{x'} = 0 \text{ se } x\neq x',
\qquad
\sum_{x\in X^{(n)}} e_x = \mathbf 1.
\]
São, portanto, idempotentes mutuamente ortogonais que somam à unidade da
álgebra.
A distinção é a seguinte: como elementos de \(\mathcal V_n\) formam uma
base do espaço vectorial; como elementos de \(\mathcal A_n(\mathbb K)\)
são idempotentes da álgebra de campos.
\end{proposicao}

\begin{proof}
Verificação ponto a ponto a partir da definição de \(e_x\).
A multiplicação em \(\mathcal A_n(\mathbb K)\) é ponto a ponto, logo
\((e_x\, e_{x'})(y)=e_x(y)\,e_{x'}(y)\), que vale \(1\) apenas quando
\(y=x=x'\) e \(0\) nos restantes casos.
\end{proof}

\begin{proposicao}[Extensão por zero ao longo da inclusão]
\label{prop:extensao-zero}
Dado \(V\in\mathcal V_n\), define-se \(\widetilde V\in\mathcal V_{n+1}\) por
\[
\widetilde V(b,\xi) :=
\begin{cases}
V(\xi) & \text{se } d(b)=0,\\
0 & \text{se } d(b)\neq 0.
\end{cases}
\]
Então \(\widetilde V\circ\iota_n = V\).
Além disso, \(\widetilde V\) anula-se fora da imagem de \(\iota_n\).
\end{proposicao}

\begin{proof}
Para \(\xi\in X^{(n)}\),
\[
\widetilde V\bigl(\iota_n(\xi)\bigr)
= \widetilde V(b_*,\xi)
= V(\xi),
\]
porque \(d(b_*)=0\).
Se \(d(b)\neq 0\), então \((b,\xi)\) não pertence à imagem de \(\iota_n\) e
\(\widetilde V(b,\xi)=0\) por definição.
\end{proof}

\begin{observacao}[Propriedades da extensão por zero]
\label{obs:distincao-extensoes}
A extensão da Proposição~\ref{prop:extensao-zero} é a \emph{extensão por zero}
fora da imagem da inclusão.
Uma \emph{extensão constante} (atribuir o mesmo valor \(V(\xi)\) a todos os
símbolos \(b\)) é uma operação distinta e não é considerada aqui.

A extensão por zero \(V\mapsto\widetilde V\) é linear.
Além disso, para o produto ponto a ponto,
\[
\widetilde{VW}=\widetilde V\,\widetilde W,
\]
porque fora da imagem de \(\iota_n\) ambos os lados se anulam e, na imagem,
coincidem com o produto em \(\mathcal V_n\).
Em contrapartida, a extensão \emph{não} preserva a unidade global: o campo
constante \(\mathbf 1\) de \(\mathcal A_n(\mathbb K)\) é enviado para o
indicador da imagem de \(\iota_n\), e não para o campo constante
\(\mathbf 1\) de \(\mathcal A_{n+1}(\mathbb K)\).
Logo a extensão por zero é um morfismo de anéis sem unidade (ou de álgebras
não unital), mas não um morfismo de álgebras com unidade.
\end{observacao}

\begin{observacao}[Valores vectoriais]
\label{obs:valores-vectoriais}
Para \(r\ge 1\) pode-se considerar \(\mathcal V_n^{(r)}=(\mathbb K^r)^{X^{(n)}}\).
A dimensão passa a ser \(r\cdot q^n\).
A estrutura do suporte e a coordenada permanecem inalteradas.
\end{observacao}

% ═══════════════════════════════════════════════════════════════════════════
\section{Álgebra de valores}
\label{sec:valores}

Seja \(K\) um corpo comutativo de característica zero.
A característica zero garante que \(k!\) é invertível em \(K\) para todo
\(k\ge 0\) (condição usada em somas parciais algébricas) e exclui a
degenerescência da característica \(2\).

\begin{definicao}[Extensão quadrática de valor]
\label{def:AmK}
Para \(m\in K\),
\[
A_m(K) := K[X]/(X^2-mX-1).
\]
Denote por \(u=[X]\) a classe da indeterminada.
\end{definicao}

\begin{proposicao}[Forma normal de \(A_m(K)\)]
\label{prop:forma-Am}
Todo elemento de \(A_m(K)\) escreve-se unicamente na forma \(a+bu\) com
\(a,b\in K\).
A multiplicação é determinada por \(u^2=mu+1\).
\end{proposicao}

\begin{proof}
O algoritmo de divisão euclidiana em \(K[X]\) produz resto de grau menor
que \(2\). A relação \(X^2-mX-1=0\) no quociente dá \(u^2=mu+1\).
\end{proof}

\begin{lema}[Teorema Chinês dos Restos para ideais principais comaximais]
\label{lema:tcr}
Seja \(R\) um anel comutativo com unidade e sejam \(a,b\in R\) tais que os
ideais \((a)\) e \((b)\) são comaximais, isto é, \((a)+(b)=R\).
Então
\[
R/(ab)\;\cong\; R/(a)\times R/(b)
\]
como anéis (e como \(R\)-álgebras).
O isomorfismo envia a classe de \(r\) no par das classes de \(r\) módulo
\((a)\) e módulo \((b)\).
\end{lema}

\begin{proof}
O morfismo \(\varphi\colon R\to R/(a)\times R/(b)\),
\(r\mapsto(r+(a),\,r+(b))\), tem núcleo \((a)\cap(b)\).
Como \((a)\) e \((b)\) são comaximais, \((a)\cap(b)=(ab)\).
Para a sobrejectividade: dados \(r_1,r_2\in R\), a comaximalidade fornece
\(x,y\in R\) com \(xa+yb=1\).
O elemento \(r:=r_2\,xa+r_1\,yb\) satisfaz
\[
r-r_1=(r_2-r_1)xa\in(a),
\qquad
r-r_2=(r_1-r_2)yb\in(b),
\]
logo \(\varphi(r)=(r_1+(a),\,r_2+(b))\).
\end{proof}

\begin{proposicao}[Decomposição quando o discriminante é quadrado]
\label{prop:decomposicao}
Suponha \(\mathrm{char}\,K\neq 2\) e \(\Delta=m^2+4\in K^{\times 2}\), isto é,
existe \(s\in K^\times\) com \(s^2=\Delta\).
Sejam \(\alpha=(m+s)/2\) e \(\beta=(m-s)/2\).
Então
\[
X^2-mX-1=(X-\alpha)(X-\beta)
\]
e o Lema~\ref{lema:tcr} fornece um isomorfismo de \(K\)-álgebras
\[
A_m(K)\xrightarrow{\cong} K\times K,
\qquad
[f]\longmapsto\bigl(f(\alpha),f(\beta)\bigr).
\]
\end{proposicao}

\begin{proof}
Como \(\mathrm{char}\,K\neq 2\), a divisão por \(2\) é legítima.
Os ideais \((X-\alpha)\) e \((X-\beta)\) em \(K[X]\) são comaximais porque
\(\alpha-\beta=s\neq 0\).
O Lema~\ref{lema:tcr} aplica-se com \(R=K[X]\).
\end{proof}

\begin{proposicao}[Idempotentes do valor]
\label{prop:idempotentes-valor}
Nas condições da Proposição~\ref{prop:decomposicao}, os elementos
\[
e_\alpha := \frac{u-\beta}{\alpha-\beta},
\qquad
e_\beta := \frac{\alpha-u}{\alpha-\beta}
\]
satisfazem
\[
e_\alpha^2=e_\alpha,\quad
e_\beta^2=e_\beta,\quad
e_\alpha e_\beta=0,\quad
e_\alpha+e_\beta=1,
\]
e \(A_m(K)=K e_\alpha\oplus K e_\beta\).
\end{proposicao}

\begin{proof}
Verificação directa usando \(u^2=mu+1\), \(\alpha+\beta=m\),
\(\alpha\beta=-1\) e \(\alpha-\beta=s\neq 0\).
\end{proof}

\begin{teorema}[Duplicação de índice]
\label{thm:duplicacao}
Nas condições da Proposição~\ref{prop:decomposicao}, existe um isomorfismo
natural de \(K\)-álgebras
\[
\mathsf T\colon A_m(K)^{X^{(n)}}
\xrightarrow{\;\cong\;}
K^{X^{(n)}\times\{0,1\}}.
\]
Explicitamente, se \(\psi\colon A_m(K)\xrightarrow{\cong}K\times K\) é o
isomorfismo da Proposição~\ref{prop:decomposicao}, então
\[
\bigl(\mathsf T(F)\bigr)(x,i) := \psi_i\bigl(F(x)\bigr).
\]
\end{teorema}

\begin{proof}
\(\mathsf T\) é morfismo de \(K\)-álgebras por definição ponto a ponto.
A inversa
\[
\bigl(\mathsf T^{-1}(G)\bigr)(x)
:= \psi^{-1}\bigl(G(x,0),G(x,1)\bigr)
\]
está bem definida e verifica as identidades de composição.
\end{proof}

\begin{observacao}[Produto cartesiano no domínio]
\label{obs:dominio}
O isomorfismo identifica \((K\times K)^{X^{(n)}}\) com
\(K^{X^{(n)}\times\{0,1\}}\).
O produto cartesiano aparece no \emph{domínio} da função após o isomorfismo,
não no contradomínio.
A definição original de \(X^{(n)}\) permanece inalterada.
\end{observacao}

\begin{proposicao}[Base de idempotentes ampliada]
\label{prop:idempotentes-ampliados}
Nas condições do Teorema~\ref{thm:duplicacao}, os elementos
\[
e_{(x,\alpha)} := e_x\, e_\alpha,
\qquad
e_{(x,\beta)} := e_x\, e_\beta
\]
(\(x\in X^{(n)}\)) formam uma família de \(2q^n\) idempotentes mutuamente
ortogonais que somam à unidade.
Em particular são linearmente independentes sobre \(K\) e
\[
\dim_K A_m(K)^{X^{(n)}} = 2q^n.
\]
\end{proposicao}

\begin{proof}
As relações de idempotência e ortogonalidade seguem das correspondentes
propriedades de \(\{e_x\}\) e \(\{e_\alpha,e_\beta\}\) por multiplicação
ponto a ponto.
Cada \(e_{(x,\alpha)}\) (resp.\ \(e_{(x,\beta)}\)) é não nulo:
no ponto \(x\), o seu valor é \(e_\alpha\) (resp.\ \(e_\beta\)).
Como \(s=\alpha-\beta\neq 0\) e todo elemento de \(A_m(K)\) tem forma
normal única \(a+bu\), com \(1\) e \(u\) linearmente independentes sobre
\(K\), as fórmulas
\[
e_\alpha=(u-\beta)/s,\qquad e_\beta=(\alpha-u)/s
\]
implicam \(e_\alpha\neq 0\) e \(e_\beta\neq 0\)
(o coeficiente de \(u\) é \(1/s\neq 0\)).
A dimensão de \(A_m(K)^{X^{(n)}}\) é \(\dim_K A_m(K)\cdot q^n=2q^n\).
Uma família de \(2q^n\) elementos mutuamente ortogonais não nulos que somam
à unidade é linearmente independente (se \(\sum\lambda_i e_i=0\), multiplicar
por \(e_j\) dá \(\lambda_j e_j=0\), logo \(\lambda_j=0\)).
\end{proof}

% --- construção de i ---

\begin{definicao}[Álgebra de valores por \(X^2+1\)]
\label{def:Kconstr}
\[
K_{\mathrm{constr}} := K[X]/(X^2+1).
\]
Denote \(i:=[X]\).
\end{definicao}

\begin{proposicao}[Relação fundamental]
\label{prop:i}
\(i^2=-1\).
\end{proposicao}

\begin{proof}
\(i^2+1=[X^2+1]=[0]\).
\end{proof}

\begin{proposicao}[Forma normal]
\label{prop:forma-normal}
Para todo \(f\in K[X]\) existem únicos \(a,b\in K\) tais que
\([f]=[a+bX]\).
\end{proposicao}

\begin{proof}
Divisão euclidiana: resto de grau menor que \(2\).
\end{proof}

\begin{proposicao}[Multiplicação nas classes]
\label{prop:mult-classes}
Para \(a,b,c,d\in K\),
\[
[a+bX]\,[c+dX] = [(ac-bd)+(ad+bc)X].
\]
\end{proposicao}

\begin{proof}
O produto dos representantes é \(ac+(ad+bc)X+bd\,X^2\).
Como \(X^2\sim -1\), tem-se \(bd\,X^2\sim -bd\).
\end{proof}

\begin{proposicao}[Estrutura de corpo]
\label{prop:K-corpo}
Se \(X^2+1\) é irredutível sobre \(K\), então \(K_{\mathrm{constr}}\) é um
corpo.
Em característica zero, \(X^2+1\) é irredutível precisamente quando \(-1\)
não é um quadrado em \(K\).
\end{proposicao}

\begin{proof}
\(K[X]\) é domínio de integridade euclidiano.
Se \(X^2+1\) é irredutível, o ideal \((X^2+1)\) é maximal e o quociente é
corpo.
Em característica zero um polinómio quadrático sem raízes é irredutível.
\end{proof}

\begin{observacao}[Independência das camadas]
\label{obs:independencia-K}
A construção
\[
K\to K[X]\to(X^2+1)\to K[X]/(X^2+1)\to a+bi
\]
pertence exclusivamente à camada dos corpos de valores.
O suporte \(X^{(n)}\) e a coordenada \(C^{(n)}\) permanecem inalterados.
\end{observacao}

\begin{lema}[Soma parcial por paridade]
\label{lema:soma-parcial}
Seja \(A\) uma álgebra associativa com unidade sobre \(K\) e seja
\(J\in A\) com \(J^2=-1_A\).
Para \(t\in K\) e \(N\ge 0\),
\[
P_N(t,J)
:= \sum_{k=0}^{N}\frac{(tJ)^k}{k!}
= \Biggl(\sum_{0\le 2k\le N}\frac{(-1)^k t^{2k}}{(2k)!}\Biggr)1_A
+ J\Biggl(\sum_{0\le 2k+1\le N}\frac{(-1)^k t^{2k+1}}{(2k+1)!}\Biggr).
\]
A identidade é puramente algébrica (soma finita); não intervêm limites nem
convergência.
\end{lema}

\begin{proof}
Particiona-se a soma pela paridade do índice.
De \(J^2=-1_A\) segue \(J^{2m}=(-1)^m 1_A\) e \(J^{2m+1}=(-1)^m J\).
Os escalares de \(K\) são centrais, logo \((tJ)^m=t^m J^m\).
\end{proof}

\begin{observacao}[Soma parcial versus exponencial]
\label{obs:parcial-vs-exp}
O Lema~\ref{lema:soma-parcial} estabelece uma identidade finita.
Não se afirma a convergência da série infinita
\(\sum_{k\ge0}(tJ)^k/k!\), nem a existência de uma exponencial em
\(\mathcal A_n(K_{\mathrm{constr}})\) sem topologia ou formalismo de
séries formais adicionais.
\end{observacao}


% ═══════════════════════════════════════════════════════════════════════════

% ═══════════════════════════════════════════════════════════════════════════
\section{Sistema indutivo, colimite e categoria restrita}
\label{sec:categoria}

\begin{definicao}[Sistema indutivo das inclusões]
\label{def:sistema-indutivo}
A sequência
\[
X^{(2)}\xrightarrow{\iota_2}X^{(3)}\xrightarrow{\iota_3}X^{(4)}\xrightarrow{\iota_4}\cdots
\]
é um sistema indutivo de conjuntos (cada \(\iota_n\) é injectiva).
\end{definicao}

\begin{proposicao}[Existência e propriedade universal do colimite de conjuntos]
\label{prop:colimite-conjuntos}
Seja
\[
U=\bigsqcup_{n\ge 2}X^{(n)}
\]
a união disjunta e seja \(\sim\) a relação de equivalência gerada por
\((n,x)\sim(n+1,\iota_n(x))\).
O quociente
\[
X^{(\infty)}
:=
U/{\sim}
=
\varinjlim_{n\ge 2}X^{(n)}
\]
é o colimite de conjuntos do sistema.
As aplicações canónicas \(\lambda_n\colon X^{(n)}\to X^{(\infty)}\) enviam
\(x\) na classe de \((n,x)\).

Além disso, se \(Z\) é um conjunto e \(f_n\colon X^{(n)}\to Z\) são funções
com \(f_{n+1}\circ\iota_n=f_n\) para todo \(n\), então existe uma única
função \(\overline f\colon X^{(\infty)}\to Z\) tal que
\(\overline f\circ\lambda_n=f_n\) para todo \(n\).
\end{proposicao}

\begin{proof}
A construção do colimite de um sistema indutivo de conjuntos é padrão.
A aplicação \(\overline f([(n,x)])=f_n(x)\) está bem definida pela condição
de compatibilidade e é a única que factoriza as \(f_n\).
\end{proof}

\begin{observacao}[O colimite não herda a coordenada dos níveis]
\label{obs:colimite-coordenada}
A existência de \(X^{(\infty)}\) como conjunto não produz, por si só, uma
coordenada
\(C^{(\infty)}\colon X^{(\infty)}\to\mathbb N\)
compatível com todas as \(C^{(n)}\) no sentido
\(C^{(\infty)}\circ\lambda_n=C^{(n)}\).
Como \(\lambda_{n+1}\circ\iota_n=\lambda_n\), uma função
\(C^{(\infty)}\colon X^{(\infty)}\to\mathbb N\) que satisfizesse
\(C^{(\infty)}\circ\lambda_n=C^{(n)}\) para todo \(n\) teria de verificar
\[
C^{(n)}(x)
=C^{(\infty)}(\lambda_n(x))
=C^{(\infty)}(\lambda_{n+1}(\iota_n(x)))
=C^{(n+1)}(\iota_n(x))
=q\,C^{(n)}(x).
\]
Como \(q\ge 2\) e \(C^{(n)}\) assume o valor \(1\) (existe \(x\in X^{(n)}\)
com \(C^{(n)}(x)=1\)), essa igualdade é impossível.
Portanto, não existe uma função no colimite que restrinja exactamente a
\(C^{(n)}\) em todos os níveis.
Isto não exclui a existência de outras coordenadas com propriedades
distintas.
\end{observacao}

\begin{definicao}[Categoria dos níveis da família]
\label{def:categoria-niveis}
Seja \(\mathbf{SC}_q^{\mathrm{fam}}\) a categoria cujos
\begin{itemize}[leftmargin=1.5em]
\item objectos são os pares \(\bigl(X^{(n)},C^{(n)}\bigr)\) para \(n\ge 2\);
\item morfismos de \(\bigl(X^{(n)},C^{(n)}\bigr)\) para
\(\bigl(X^{(m)},C^{(m)}\bigr)\) (\(m\ge n\)) são as composições de
inclusões
\[
\iota_{m-1}\circ\cdots\circ\iota_n
\]
(quando \(m>n\)) e a identidade (quando \(m=n\));
\item se \(m<n\), não há morfismos.
\end{itemize}
A composição é a composição usual de funções; as identidades são as
funções identidade.
\end{definicao}

\begin{proposicao}[A estrutura é uma categoria]
\label{prop:categoria}
\(\mathbf{SC}_q^{\mathrm{fam}}\) é uma categoria.
\end{proposicao}

\begin{proof}
A composição de composições de inclusões é novamente uma composição de
inclusões (associatividade da composição de funções).
As identidades satisfazem as leis unitárias por definição.
\end{proof}

\begin{observacao}[Âmbito da categoria]
\label{obs:ambito-categoria}
A categoria \(\mathbf{SC}_q^{\mathrm{fam}}\) contém apenas os níveis da
família recursiva e as inclusões canónicas de peso mínimo.
Não se afirma que ela esgote todos os morfismos possíveis entre suportes
coordenados arbitrários.
\end{observacao}

% ═══════════════════════════════════════════════════════════════════════════
\section{Realizações e campos de contagem}
\label{sec:realizacoes}

\begin{definicao}[Realização]
\label{def:realizacao}
Uma \emph{realização} de comprimento \(N\) sobre \(X^{(n)}\) é uma função
\[
\pi\colon\{1,\dots,N\}\to X^{(n)}.
\]
\end{definicao}

\begin{definicao}[Campo de contagem]
\label{def:campo-contagem}
O campo de contagem associado a \(\pi\) é
\[
G_\pi\colon X^{(n)}\to\mathbb N_0,
\qquad
G_\pi(x) := \lvert\pi^{-1}(x)\rvert.
\]
Assim \(G_\pi\) define um elemento de \(\mathcal A_n(\mathbb Z)\) sob a
inclusão natural \(\mathbb N_0\subset\mathbb Z\).
\end{definicao}

\begin{proposicao}[Soma de realizações]
\label{prop:soma-realizacoes}
Se \(\pi\) e \(\rho\) são realizações (de comprimentos \(N\) e \(M\)) sobre
o mesmo suporte, a concatenação \(\pi\ast\rho\) satisfaz
\[
G_{\pi\ast\rho} = G_\pi + G_\rho
\]
(ponto a ponto).
\end{proposicao}

\begin{proof}
A pré-imagem de \(x\) sob a concatenação é a união disjunta das pré-imagens
sob \(\pi\) e sob \(\rho\).
\end{proof}

% ═══════════════════════════════════════════════════════════════════════════
\section{Questões em aberto}
\label{sec:abertas}

\begin{questao}[Estrutura e coordenada no colimite]
\label{questao:colimite}
Qual estrutura (se alguma) e qual coordenada (se alguma) sobrevivem de forma
canónica no colimite \(X^{(\infty)}\)?
A existência do colimite de conjuntos está estabelecida
(Proposição~\ref{prop:colimite-conjuntos}); a questão interessante é a
preservação de estrutura adicional.
\end{questao}

\begin{questao}[Categoria mais ampla]
\label{questao:categoria-ampla}
É possível estender \(\mathbf{SC}_q^{\mathrm{fam}}\) a uma categoria que
inclua morfismos entre suportes coordenados não relacionados pela família
recursiva, de modo que os morfismos de escala da
Definição~\ref{def:morfismos} sejam morfismos dessa categoria?
\end{questao}

\begin{questao}[Extensão da soma parcial]
\label{questao:analitica}
A soma parcial do Lema~\ref{lema:soma-parcial} pode ser completada a uma
exponencial formal em \(\mathcal A_n(K_{\mathrm{constr}})\) de forma
rigorosa sem introduzir topologia adicional?
\end{questao}

\section*{Nota final}

Os resultados das Secções~\ref{sec:axiomas}--\ref{sec:realizacoes}
decorrem das definições e hipóteses enunciadas no próprio texto.
As questões da Secção~\ref{sec:abertas} delimitam o que permanece em aberto
sem invocar estruturas externas ao artigo.

\end{document}
